\scanpage{70}
\uline{Morphismes ess. \deleted{\unsure{affines}} de prés. finie}

Soit $f:X\longrightarrow Y$ un morphisme. Considérons les conditions

(i) $f$ est la $\varprojlim$ d'un syst. projectif filtrant ess. affine de $Y$-préschémas de \uline{prés. finie}

(ii) $f$ est égal à un composé $X\xrightarrow{f_1}X_0\xrightarrow{f_0}Y$, avec $f_0$ de prés. finie et $f_1$ affine

(iii) Comme (ii), mais avec $X=\operatorname{Spec}(\mathcal A)$, $\mathcal A$ Alg. quasi-cohérente sur $X_0$ qui est $\varinjlim$ d'une famille filtrante d'Algèbres de prés. finie sur $X_0$.

Alors $\text{(i)}\Leftrightarrow\text{(iii)}\Longrightarrow\text{(ii)}$, et tout est équivalent si $Y$ est loc. noeth.

\uline{Définition} On dit que $f$ est ess. de prés. finie s'il satisfait la condition (i). On dit que $X$ est ess. de prés. finie si $X\longrightarrow\operatorname{Spec}\mathbb{Z}$ est ess. de pr. finie.

N.B. \deleted{Alors} $f$ ess. \unsure{pr. fo|pr. f.} $\Rightarrow f$ qu. cpct qu. sép.

\uline{Prop} (i) Prés finie $\Rightarrow$ ess. prés. finie ; affine + but loc. noeth ou qu. cpt \inserted{qu. sép} avec faisceau inv. ample $\Longrightarrow$ ess. pr. f. ; \deleted{\unsure{$f$}} imm. \unsure{qu. cpcte.|\deleted{qu. cpte}} + but qu. cpt qu. sép. $\Longrightarrow$ ess. prés. finie.

(ii) Stabilité chang$^{\mathrm{mt}}$ de base

(iii) Stabilité composition

(iv) Si $g\deleted{\unsure{f}}:Y\mathrel{\unsure{-|\to}}S$ ess. pr. f. \inserted{et \unsure{$S$} qu. cpt qu. sép., ou si $g$ est de pr. f.}, alors : $gf$ \deleted{\unsure{$f$|est}} est \inserted{ess. pr. f.} ssi $f$ l'est

\uline{Prop} Supposons $S$ \unsure{univ.} japonais, p. ex. $S=\operatorname{Spec}\mathbb Z$. Alors si $X$ est \deleted{géom. unibr.} ess. pr. f. sur $S$, pour que $X$ soit géom unibr, il f. et s. qu'il soit de la forme $\varprojlim X_i$, avec les conditions habituelles, et de plus les $X_i$ géom. unibranches (resp. normaux).
% archive sheet number: 69
